\documentclass[12pt]{article}
\usepackage[utf8]{inputenc}
\usepackage[T1]{fontenc}
\usepackage{hyperref}
\title{Tereweng Language Resources for K-12 Curriculum Development}
\author{Compiled by Claude Sonnet 5 (claude-sonnet-5) for Andr\'as Kornai}
\date{}
\begin{document}
\maketitle

\section{Introduction}
Tereweng (ISO 639-3 twg) is an Alor--Pantar language spoken on Tereweng Island, southeast of Pantar, in Indonesia. It is closely related to Blagar and is sometimes treated as a dialect of it. Speaker estimates vary (about 800 to 1,200). Most speakers also use Alor Malay and Indonesian. Searches found no school materials in or for Tereweng. This report is therefore short.

\section{Primary Educational Resources}
No primers, textbooks or school programmes for Tereweng were found in search results.

\section{Linguistic and Reference Materials}
\begin{itemize}
\item Wikipedia, Tereweng language (overview, classification, relation to Blagar): \url{https://en.wikipedia.org/wiki/Tereweng_language}
\item Wikipedia, Alor--Pantar languages (family context): \url{https://en.wikipedia.org/wiki/Alor%E2%80%93Pantar_languages}
\item ASJP, Blagar Tereweng (a word list entry): \url{https://asjp.clld.org/languages/BLAGAR_TEREWENG}
\end{itemize}

\section{Religious and Community Resources}
Joshua Project has a Tereweng page, but it places the group on Gili Trawangan, which looks wrong, so treat it with caution: \url{https://joshuaproject.net/people_groups/15443}. No scripture or community organisation was found.

\section{Digital and Audio Resources}
None found.

\section{Curriculum Development Context}
There is almost no documented material. A curriculum would need new fieldwork with the community, possibly building on Blagar documentation. Indonesian is the likely language of instruction in the meantime.

\end{document}